\documentclass[11pt,a4paper]{article}
\usepackage[margin=2.3cm]{geometry}
\usepackage{fontspec}
\usepackage{booktabs}
\usepackage{amsmath}
\usepackage[colorlinks=true,linkcolor=blue!50!black,urlcolor=blue!50!black]{hyperref}

\setmainfont{Thonburi}
\XeTeXlinebreaklocale "th"
\XeTeXlinebreakskip = 0pt plus 1pt

\title{\textbf{พลังของ Prompt: จาก Zero-shot สู่ Few-shot}\\[4pt]
\large บทความเชิงวิเคราะห์เพื่อการเรียนรู้การใช้ AI อย่างมีประสิทธิภาพ}
\author{อิทธิพล วงศ์อภัย \\ \small รายวิชา GenAI 1 — สัปดาห์ที่ 4}
\date{8 กันยายน 2569}

\begin{document}
\maketitle

\begin{abstract}
บทความนี้วิเคราะห์ความแตกต่างระหว่างการเขียนพรอมป์แบบ \emph{zero-shot} (สั่งงานครั้งเดียวโดยไม่มีตัวอย่าง)
กับแบบ \emph{few-shot} (ใส่บริบท บทบาท และตัวอย่างประกอบ) โดยใช้ผลการทดลองจริงจากการบ้าน 6 แง่
ของรายวิชา พร้อมทั้งเสนอกรอบการปรับปรุงพรอมป์อย่างเป็นระบบ
\end{abstract}

\section{บทนำ}
ถ้อยคำเดียวกันอาจให้ผลลัพธ์ต่างกันอย่างมีนัยสำคัญ เมื่อบริบทของคำสั่งเปลี่ยน
งานเขียนชิ้นนี้สรุปประสบการณ์จากการทดลองเขียนพรอมป์สองรอบในหกเครื่องมือ
ได้แก่ การสร้างภาพ กราฟคณิตศาสตร์ ไดอะแกรม เอกสาร LaTeX สไลด์สรุป และเว็บไซต์

\section{กรอบการวิเคราะห์}
กำหนดให้คุณภาพของผลลัพธ์รอบที่ $n$ เป็นฟังก์ชันของปริมาณบริบทที่ใส่เข้าไป
\begin{equation}
Q_n \;=\; Q_0 \,(1+\alpha)^{n}, \qquad \alpha = \frac{\Delta c}{C + \Delta c}
\end{equation}
เมื่อ $Q_0$ คือคุณภาพเริ่มต้นจากพรอมป์แรก $\Delta c$ คือบริบทที่เพิ่มในแต่ละรอบ
และ $C$ คือค่าคงที่ของโมเดล สอดคล้องกับเส้นโค้งการเรียนรู้ที่เพิ่มขึ้นแบบ asymptotic
ต่อเป้าหมายที่แท้จริงของผู้ใช้

\section{ผลการทดลอง}
\begin{table}[h]
\centering
\begin{tabular}{@{}llll@{}}
\toprule
\textbf{แง่} & \textbf{Zero-shot} & \textbf{Few-shot} & \textbf{การปรับที่สำคัญ} \\
\midrule
ภาพ AI & ฉากทั่วไป ตัวอักษรเพี้ยน & ฉากตรงคอนเซปต์ คมชัด & ระบุแสง มุมกล้อง องค์ประกอบ \\
กราฟ Desmos & เส้นเดียว อ่านยาก & กราฟเชิงโต้ตอบ มี slider & กำหนดช่วง สี และพารามิเตอร์ \\
ไดอะแกรม & กล่อง 4 กล่องธรรมดา & flowchart 2 รอบ มีสี & ระบุโครงสร้างและ style \\
LaTeX & ข้อความสั้น ภาษาเดียว & บทความครบ ไทย+คณิต & กำหนดโครงเรื่องและฟอนต์ \\
สไลด์สรุป & หัวข้อกระจัดกระจาย & สไลด์เรียงตามเนื้อหา & ระบุแหล่งที่มาและรูปแบบ \\
เว็บไซต์ & หน้าเดียว ไม่มีเมนู & 6 หมวด มีเมนูนำทาง & ระบุสถาปัตยกรรมหน้าเว็บ \\
\bottomrule
\end{tabular}
\caption{เปรียบเทียบผลลัพธ์รอบแรกกับรอบปรับปรุงในทั้ง 6 แง่}
\end{table}

\section{ข้อเสนอแนะ}
\begin{itemize}
  \item เริ่มด้วยพรอมป์สั้นเพื่อสำรวจขอบเขต แล้วปรับปรุงเชิงลึกไม่เกิน 2--3 รอบ
  \item ใส่บทบาท กลุ่มเป้าหมาย รูปแบบผลลัพธ์ และข้อจำกัด ทุกครั้งที่สั่งงานเชิงซับซ้อน
  \item บันทึกพรอมป์ที่ได้ผลไว้เป็นคลังตัวอย่างสำหรับงานอื่นในอนาคต
\end{itemize}

\section{บทสรุป}
การเพิ่มบริบทอย่างมีวิธีทำให้คุณภาพผลลัพธ์เพิ่มขึ้นตามสมการ (1)
แต่ต้นทุนเวลาที่เพิ่มขึ้นเร็วกว่าผลตอบแทนชี้ว่า จุดที่เหมาะสมของการวนปรับพรอมป์
อยู่ที่รอบที่สองถึงสามเป็นส่วนใหญ่

\end{document}
